\section{Results}
\label{sec:results}

\subsection{Live-Edge Results}
\label{sec:sw_results}

\autoref{tab:sw} reports live-edge and settled errors for the three
flights using the configuration in \autoref{sec:setup}.
\autoref{fig:traj} and \autoref{fig:error_time} show the trajectories
and live-edge errors.  We report velocity first because Doppler
constrains it directly.  Position is its integral.

\begin{table}[t]
\centering
\caption{Sliding Window: Live Leading-Edge and Settled \ac{rmse}
  (per-window solve time on a laptop CPU)}
\label{tab:sw}
\small
\setlength{\tabcolsep}{3pt}
\begin{tabular}{@{}llcccc@{}}
\toprule
Flight & Mode & Vel (m/s) & Ori (\degree) & Drift (\%) & $t_{\mathrm{win}}$ (s) \\
\midrule
\multirow{2}{*}{Slow racing}
  & \textbf{Live edge} & \textbf{0.16} & \textbf{1.65} & \textbf{0.32} & \multirow{2}{*}{0.16} \\
  & Settled & 0.17\textsuperscript{$\ast$} & 1.16 & 0.29 & \\
\midrule
\multirow{2}{*}{Fast racing}
  & \textbf{Live edge} & \textbf{0.28} & \textbf{2.56} & \textbf{1.33} & \multirow{2}{*}{0.18} \\
  & Settled & 0.42\textsuperscript{$\ast$} & 2.41 & 1.28 & \\
\midrule
\multirow{2}{*}{Backflips}
  & \textbf{Live edge} & \textbf{1.66} & \textbf{12.1} & \textbf{2.91} & \multirow{2}{*}{0.20} \\
  & Settled & 1.72\textsuperscript{$\ast$} & 7.91 & 2.75 & \\
\bottomrule
\multicolumn{6}{@{}p{\columnwidth}@{}}{\scriptsize
  The \textbf{bold} live-edge rows are the deployment-relevant figures.
  $^\ast$The settled trajectory combines control points last optimized
  in different windows.  Disagreement between those solutions distorts
  its velocity near stride boundaries and inflates the settled velocity
  \ac{rmse}, while live velocity comes from each window's own spline.
  Path lengths 48.1 / 44.9 / 54.1\,m.  Repeated solver runs on the same
  recordings vary by $<$2\%.}
\end{tabular}
\end{table}

On the racing flights, velocity \ac{rmse} is $0.16$--$0.28$\,m/s and
position drift is $0.3$--$1.3\%$ of the path length.  Errors increase
on the backflip flight, with orientation \ac{rmse} reaching $12.1\degree$
(\autoref{sec:backflips}).  All three regimes solve faster than the
0.3\,s stride on a laptop CPU (\autoref{tab:sw}), i.e.\ the deployed
settings run in real time.

\begin{figure*}[t]
  \centering
  \includegraphics[width=\textwidth]{figures/traj_combined.pdf}
  \caption{Trajectory overlays for all three flights (columns): \ac{mocap}
    ground truth vs.\ the sliding-window live edge, X--Y (top) and Y--Z
    (bottom) projections. Within each column both projections share a
    common $Y$ range and metric scale (equal aspect, so loops keep their
    true shape).  The scale differs between columns.  Start~$\blacksquare$.
    the dotted lead-in is the first window's estimate, available only
    once that window has been solved (the smoother needs one complete
    3\,s window), so it is not a live output at those timestamps.  It
    starts on the ground truth by the start-anchored alignment.
    The backflip flight circles in the horizontal plane while looping
    in the vertical plane.  Its position drifts through the flips
    ($2.9\%$ of the path) while the loops remain visible in the
    orientation estimate (\autoref{sec:backflips}).}
  \label{fig:traj}
\end{figure*}

\begin{figure*}[t]
  \centering
  \includegraphics[width=\textwidth]{figures/error_over_time.pdf}
  \caption{Live-edge velocity, orientation, and position error (the last in percent of the flight's path length, the unit of the drift metric)
    vs.\ time for slow racing (top), fast racing (middle), and backflips
    (bottom).  Velocity and orientation errors spike at high-speed,
    high-dynamics sections.  Fast-racing position error accumulates with
    trajectory length.  Backflip errors grow step-wise \emph{during} the
    flips, where the radar measurement model degrades
    (\autoref{sec:backflips}).}
  \label{fig:error_time}
\end{figure*}

\subsection{Held-Out Validation}
\label{sec:heldout}
The per-return law and the stream weights were fitted on the two
racing flights with the backflip flight held out.  The frame-shared
term, the alias floor, the bias-prior widths and the staleness rates
were measured on all three.  The configuration was then run
\emph{unmodified} on flights never used for any decision.

A second, independent fast-racing flight was recorded with the radar's
on-chip clustering disabled (61 returns per frame instead of 11).  The
front end then collapses duplicate detections to the physical cluster
count, a preprocessing step that applies only when on-chip clustering
is off, with the solver configuration unchanged.  Velocity \ac{rmse}
is $0.27$\,m/s, compared with $0.28$\,m/s on the fitting flight.
Orientation \ac{rmse} is $2.7\degree$ and drift is $1.2\%$.
The four public ICINS flights in \autoref{sec:baselines} test
transfer to a different platform.

Accuracy degrades on a held-out circular flight at 5\,m/s.
With the radar's Doppler ambiguity of 3.1\,m/s, $63.8\%$ of its
returns must be
alias-unwrapped, the \ac{imu}-integrated unwrapper's alias decisions
degrade, and the live edge reaches
$1.30$\,m/s\,/\,$11.7\degree$\,/\,$9.1\%$.
Down-weighting aliased returns cannot compensate for incorrect wrap selection.



\subsection{External Cross-Validation Against EKF Baselines}
\label{sec:baselines}
We cross-validate against the \ac{ekf}-\ac{rio} family of Doer and
Trommer~\cite{doer2020ekf,doer2021yrio} in both directions on shared
data, with their toolbox pinned at a fixed commit and every parameter
deviation documented in the repository.

\textbf{Our system on their data.}  We run our solver on the four public ICINS-2021 indoor flights~\cite{doer2021yrio} (the same radar chip, a different \ac{imu}
at 409\,Hz, VINS pseudo ground truth). We use their published extrinsic
calibration, \emph{all other parameters remain those measured on our sensor and flights}; their own frame-shared term is $\sigma_c = 0.93$, and we kept
$1.27$.  Their
estimators (\mbox{ekf-rio} with unaided yaw and \mbox{ekf-yrio} with
radar yaw aiding) are re-run from their stock configurations and evaluated with
the same causal protocol.  \autoref{tab:baselines} shows comparable velocity errors on every flight.
Drift is $0.4$--$1.3\%$ of the path, ahead of the unaided baseline on
flight~1 and behind both baselines on the other three, with the
difference in the vertical channel (\autoref{tab:baselines}(b),
\autoref{sec:vertical}).  The
orientation column prefers our solver only because we anchor heading 
 with a pseudo-magnetometer while theirs is sensor-only or
unaided. On the heading-independent roll and pitch the three systems are
similar (\autoref{tab:baselines}(b)), and with our heading aiding
removed our yaw drifts similarly to their unaided filter
(\autoref{sec:yaw}).

The consensus prefilter reduces our vertical drift on all four public
flights.  On flight~1 it falls from $14.3\%$ to $0.35\%$ of the path,
an error the baselines do not show.

\textbf{Their filters on our data.}
Ported unmodified onto our data, their filters, expecting order-100 clouds, fail on our sparse radar frames. On slow racing the $\chi^2(3)$
 gate of ekf-rio rejects $99.7\%$ of the radar ego-velocity
updates and reverts to \ac{imu} dead-reckoning. 

\begin{table}[t]
\caption{Cross-Validation on the Public ICINS-2021 Flight
Datasets~\cite{doer2021yrio}, All Systems Under Our Causal Protocol}
\label{tab:baselines}
\centering
\footnotesize
\setlength{\tabcolsep}{3pt}%
\begin{tabular*}{\columnwidth}{@{\extracolsep{\fill}}lccc}
\toprule
\multicolumn{4}{@{}l}{\textit{(a)~velocity [m/s] / orientation [\degree] / drift [\% path]}}\\
\midrule
Flight (path) & ekf-rio & ekf-yrio & Ours \\
\midrule
1 (142\,m) & 0.20 / 10.1 / 1.0 & \textbf{0.08} / 0.70 / \textbf{0.3} & 0.09 / \textbf{0.55} / 0.4 \\
2 (45\,m)  & \textbf{0.08} / 1.59 / 0.6 & \textbf{0.08} / \textbf{0.91} / \textbf{0.4} & \textbf{0.08} / 1.00 / 1.0 \\
3 (147\,m) & \textbf{0.08} / 1.07 / \textbf{0.3} & \textbf{0.08} / 1.09 / \textbf{0.3} & \textbf{0.08} / \textbf{0.53} / 1.3 \\
4 (78\,m)  & 0.09 / 3.67 / 0.5 & \textbf{0.07} / 1.41 / \textbf{0.3} & 0.08 / \textbf{1.21} / 1.2 \\
\bottomrule
\end{tabular*}

\vspace{4pt}

{\scriptsize
\setlength{\tabcolsep}{1.5pt}%
\begin{tabular*}{\columnwidth}{@{\extracolsep{\fill}}lcccccccc}
\toprule
\multicolumn{9}{@{}l}{\textit{(b)~Decomposition}\quad(systems: ekf-rio / ekf-yrio / Ours)}\\
\midrule
 & \multicolumn{3}{c}{Roll/pitch [\degree]} & \multicolumn{2}{c}{Our drift [\%]} & \multicolumn{3}{c}{Whole ATE [m]} \\
\cmidrule(lr){2-4}\cmidrule(lr){5-6}\cmidrule(lr){7-9}
Fl.\ & rio & yrio & Ours & hor.\ & ver.\ & rio & yrio & Ours \\
\midrule
1 & 0.24 & 0.25 & 0.22 & 0.21 & 0.35 & 0.87 & 0.23 & 0.41 \\
2 & 0.57 & 0.59 & 0.58 & 0.55 & 0.85 & 0.13 & 0.10 & 0.22 \\
3 & 0.25 & 0.26 & 0.24 & 0.49 & 1.28 & 0.36 & 0.26 & 0.90 \\
4 & 0.74 & 0.74 & 0.75 & 0.61 & 1.08 & 0.30 & 0.15 & 0.49 \\
\bottomrule
\end{tabular*}}

\vspace{2pt}

\parbox{\columnwidth}{\raggedright\scriptsize Live estimate, start-anchored
  alignment.  \textbf{(a)}~headline, with heading aiding differing by design (see
  text).  \textbf{(b)}~decomposition: \emph{roll/pitch} is the
  heading-independent attitude \ac{rmse} (yaw excluded), so the orientation differences in~(a) mainly reflect yaw.  \emph{Our drift} splits position into
  horizontal/vertical.  \emph{Whole-traj.\ \ac{ate}} re-evaluates every
  system under EVO/KITTI Umeyama alignment (no scale) for cross-checking
  published figures.  Prefilter \emph{disabled}, flight~1: vertical
  drift $14.3\%$, whole-traj.\ \ac{ate} $11.7$\,m.\par}
\end{table}

\subsection{Output Covariance Consistency (NEES)}
\label{sec:nees}
The reported live-edge covariance of $(\vv_w, \vR)$ is tested with the
\ac{nees}~\cite{barshalom2001estimation}.  Each window's error against
ground truth is normalized by that covariance.  For a consistent
estimator each 3-DoF block has mean 3.  On both racing
flights the reported covariance is \emph{conservative}: the velocity
\ac{nees} mean is $0.20$ on slow racing ($N{=}49$) and $0.43$ on fast
racing ($N{=}20$), the orientation mean $0.39$ and $1.04$, against
the expected 3 (nominal 95\% $\chi^2$ intervals for the $N$-window mean,
$[2.4,3.7]$ and $[2.0,4.2]$, which treat the overlapping windows as
independent).  The reported standard deviations are thus
$1.7$--$3.8\times$ too large.  On backflips only 6 of 49 windows
remain after the ground-truth mask.  Their covariance is overconfident
($1.4\times$ in velocity and $3.2\times$ in orientation, in standard
deviation), but the sample is too small for a firm conclusion.

\subsection{Ablations}
\label{sec:ablations}
Selected front-end, weighting and prior choices are reverted one at a
time at the deployed operating point (live edge: velocity /
orientation / drift).

\textbf{The noise floor $\sigma_0$.}  We use the noise floor derived
from the racing fits on all flights (\autoref{sec:char_rate}).
\autoref{tab:s0ladder} tests values of $0.5$ and $2.5$ times this
setting.  On the racing flights, velocity error changes by about 7\%
($0.165$ to $0.177$\,m/s on slow racing).  On backflips, increasing
the floor to $0.207$\,m/s reduces drift from $2.9$ to $2.0\%$, with
little change in orientation error, motivating the in-flight
noise-floor estimate listed as future work.

\begin{table}[t]
\centering
\caption{Sensitivity to the law's noise floor $\sigma_0$ (live edge:
  velocity / orientation / drift)}
\label{tab:s0ladder}
\footnotesize
\setlength{\tabcolsep}{3pt}
\resizebox{\columnwidth}{!}{%
\begin{tabular}{@{}lccc@{}}
\toprule
$\sigma_0$ & Slow racing & Fast racing & Backflips \\
\midrule
$0.5\times$ (0.041\,m/s) & 0.18 / 1.65 / 0.48 & 0.29 / 2.55 / 1.40 & 1.74 / 12.2 / 5.1 \\
\textbf{$1\times$ (0.083, deployed)} & \textbf{0.16 / 1.65 / 0.32} & \textbf{0.28 / 2.56 / 1.33} & \textbf{1.66 / 12.1 / 2.9} \\
$2.5\times$ (0.207\,m/s) & 0.16 / 1.67 / 0.30 & 0.29 / 2.58 / 1.40 & 1.65 / 12.2 / 2.0 \\
\bottomrule
\end{tabular}%
}
\end{table}

\textbf{Front end.}  \autoref{tab:frontend} reverts the front-end
elements one at a time.  The consensus prefilter is needed on the
racing flights (velocity error $+21\%$ and $+51\%$ without it) and on the
public dataset (\autoref{sec:baselines}), where the vertical channel
fails without it.  On backflips, disabling the prefilter reduces
drift from $2.9\%$ to $2.4\%$.  The filter removes returns from already
sparse backflip frames.  We accept this drift penalty to retain its
benefits on racing and the public dataset.
Robust kernels (radar $\delta = 1$\,m/s,
accelerometer $\delta_a = 2$\,m/s$^2$) have negligible effect on slow
racing.  On fast racing, they slightly worsen velocity and orientation
errors while reducing drift by $0.2$\,pp.  With both enabled,
backflip orientation error falls to $7.4\degree$ while drift rises
by $2$\,pp.  Enabling only the accelerometer kernel gives a similar
orientation error of $7.6\degree$, so it accounts for most of the
orientation improvement (\autoref{sec:backflips}).  For the
aliased returns, full weighting increases fast-racing errors and
deletion costs fast-racing drift, so the measured down-weight is kept.

\textbf{Priors.}  Replacing the measured bias priors with hand-set
values increases slow-racing velocity error by 5\% and backflip
orientation error by $1.4\degree$, but reduces fast-racing velocity
error by 4\%.  We retain the measurement-based priors because their
widths reflect observed bias variability.  The position-row staleness
has negligible effect on racing and reduces the backflip orientation error by
$1.2\degree$.  Re-applying the window-start priors at every stride
cuts the fast-racing drift from $1.33$ to $0.92\%$ by holding the
vertical velocity in place, but doubles it on slow racing and on the
flip flight and double-counts the previous estimate (\autoref{sec:sw}).

\textbf{Regularization.}  Removing the smoothness priors substantially
improves backflip velocity and position accuracy, with a smaller
improvement in orientation (velocity error $1.66$ to $0.76$\,m/s,
orientation $12.1$ to $10.5\degree$, drift $2.9$ to $1.1\%$), but
increases fast-racing position drift to $5.7\%$, predominantly in the
vertical channel.  The shared configuration therefore retains them
(\autoref{sec:fixed_constants}).

\begin{table}[t]
\centering
\caption{Front-end and prior ablation (live edge: velocity / orientation / drift)}
\label{tab:frontend}
\footnotesize
\setlength{\tabcolsep}{3pt}
\resizebox{\columnwidth}{!}{%
\begin{tabular}{@{}lccc@{}}
\toprule
Configuration & Slow racing & Fast racing & Backflips \\
\midrule
\textbf{Deployed} & \textbf{0.16 / 1.65 / 0.32} & \textbf{0.28 / 2.56 / 1.33} & \textbf{1.66 / 12.1 / 2.9} \\
No prefilter & 0.20 / 1.66 / 0.88 & 0.42 / 2.63 / 1.70 & 1.61 / 12.1 / 2.4 \\
Robust kernels on ($\delta{=}1$, $\delta_a{=}2$) & 0.16 / 1.65 / 0.32 & 0.29 / 2.69 / 1.13 & 1.96 / 7.38 / 4.9 \\
Aliased kept ($w{=}1$) & 0.17 / 1.67 / 0.27 & 0.37 / 2.86 / 2.17 & 1.61 / 12.1 / 2.9 \\
Aliased deleted & 0.16 / 1.65 / 0.35 & 0.29 / 2.55 / 1.65 & 1.66 / 12.0 / 3.1 \\
\midrule
Bias priors hand-set ($10^4$) & 0.17 / 1.65 / 0.36 & 0.27 / 2.58 / 1.25 & 1.68 / 13.5 / 2.9 \\
Position staleness off & 0.16 / 1.64 / 0.32 & 0.28 / 2.50 / 1.35 & 1.74 / 13.3 / 2.8 \\
Window-start priors every stride & 0.16 / 1.69 / 0.60 & 0.28 / 2.41 / 0.92 & 1.69 / 12.5 / 5.2 \\

Smoothness priors off & 0.16 / 1.67 / 0.34 & 0.23 / 2.52 / 5.72 & 0.76 / 10.5 / 1.1 \\
\bottomrule
\end{tabular}%
}
\end{table}
